\documentclass[10pt,a4paper]{amsart}
\usepackage[margin=19mm]{geometry}
\usepackage{amsmath,amssymb,mathtools}
\usepackage[scr=rsfs,cal=euler]{mathalfa}
\usepackage{xfrac}
\usepackage[unicode,hidelinks,bookmarks=false]{hyperref}
\newcommand{\ie}{i.e.\ }
\newcommand{\cf}{cf.\ }
\newcommand{\resp}{respectively\ }
\newcommand{\Subsection}{\subsection}
\providecommand{\nobreakdash}{\nobreak\hbox{-}}
\setlength{\emergencystretch}{2em}
\multlinegap0pt
\usepackage{mathtools}
\usepackage{booktabs,longtable,array}
\hypersetup{colorlinks=true,linkcolor=blue,citecolor=blue,urlcolor=blue,pdftitle={Separable integer partition classes and Slater's list -- II}}
\newtheorem{theorem}{Theorem}[section]
\newtheorem{lemma}[theorem]{Lemma}
\newtheorem{cor}[theorem]{Corollary}
\newtheorem{prop}[theorem]{Proposition}
\newtheorem{definition}[theorem]{Definition}
\newtheorem{example}[theorem]{Example}
\newtheorem{remark}[theorem]{Remark}
\numberwithin{equation}{section}
\allowdisplaybreaks
\newcommand\numberthis{\addtocounter{equation}{1}\tag{\theequation}}
\begin{document}
\title[Separable integer partition classes and Slater's list -- II]{Separable integer partition classes and Slater's list -- II}
\author{Aritram Dhar, Ankush Goswami, Runqiao Li}
\date{}
\maketitle
\noindent\emph{Source and licence.} Separable integer partition classes and Slater's list -- II. Version arXiv:2606.27702v2 (CC BY 4.0). This material is distributed under the Creative Commons Attribution 4.0 International licence, http://creativecommons.org/licenses/by/4.0/, as recorded in the arXiv OAI licence metadata for this exact version and shown on the abstract page https://arxiv.org/abs/2606.27702v2. Full rights notice: licenses/arxiv-2606.27702-CC-BY-4.0.txt.\par\smallskip
\noindent\textit{Editorial scope.} Counted unit: the original \texttt{theorem} environment \texttt{thm:S1-closed} at source lines 420--429 together with its complete proof at lines 431--451; the delivered document keeps source lines 363--452 so that every referenced label and every citation resolves inside it. Native editable source is the paper's own concatenated TeX; the AMS wrapper is editorial formatting only. No publisher class, upstream script or unrelated artwork is redistributed.
\section{Even-position overpartitions and identities (28), (50), and (51)}\label{sec:S1}
In this section, we treat Equation $(28)$, $(50)$, and $(51)$ in Slater's list. They are given by
\begin{equation}\label{eq:Slater28}
\sum_{n=0}^{\infty}\frac{(-q^2;q^2)_{n}q^{n(n+1)}}{(q;q)_{2n+1}}
=\dfrac{(-q^2;q^2)_{\infty}(-q,-q^5,q^6;q^6)_{\infty}}{(q^2;q^2)_{\infty}},
\end{equation}
\begin{equation}\label{eq:Slater50}
\sum_{n=0}^{\infty}\frac{(-q;q^2)_nq^{n(n+2)}}{(q;q)_{2n+1}}
=\frac{(q^2,q^{10},q^{12};q^{12})_{\infty}}{(q;q)_{\infty}},
\end{equation}
and
\begin{equation}\label{eq:Slater51}
\sum_{n=0}^{\infty}\frac{(-q;q^2)_nq^{n(n+1)}}{(q;q)_{2n+1}}
=\frac{(q^4,q^8,q^{12};q^{12})_{\infty}}{(q;q)_{\infty}}.
\end{equation}

Let $\mathcal{S}_1$ be the set of overpartitions $\lambda=(\lambda_1,\lambda_2,\ldots,\lambda_{\ell})$ such that
\begin{enumerate}
    \item Only even-indexed parts may be overlined.
    \item $\lambda_{2i}-\lambda_{2i+1}\geq1$, and $\lambda_{2i}-\lambda_{2i+1}\geq2$ if $\lambda_{2i}$ is overlined.
\end{enumerate}
The class $\mathcal{S}_1$ is an SIP class of modulus $1$, and the basis $\mathcal{B}_{\mathcal{S}_1}$ is the set of partitions in $\mathcal{S}_1$ such that
\begin{enumerate}
    \item $\lambda_{2i-1}=\lambda_{2i}$ if $\lambda_{2i}>0$.
    \item $\lambda_{2i}-\lambda_{2i+1}=2$ if $\lambda_{2i}$ is overlined, and $\lambda_{2i}-\lambda_{2i+1}=1$ otherwise.
\end{enumerate}
Let $B_{\mathcal S_1}(n,h):=B_{\mathcal S_1}(n,h;a,z,q)$ be the generating function for partitions in $\mathcal{B}_{\mathcal{S}_1}$ with length $n$ and largest part $h$, where $a$, $z$, and $q$ keep track of the number of overlined parts, the number of even-indexed parts, and the weight, respectively. We will show that
$$\sum_{\lambda\in\mathcal{S}_1}a^{\overline{\#}(\lambda)}z^{\lfloor\frac{\#(\lambda)}{2}\rfloor}q^{|\lambda|}=\sum_{n=0}^{\infty}\frac{(-aq^2;q^2)_{n}z^nq^{n(n+1)}}{(q;q)_{2n+1}}.$$
Let $(a,z)\to(1,1)$, $(a,z)\to(q^{-1},1)$, and $(a,z)\to(q^{-1},q)$, we will have the sum sides of \eqref{eq:Slater28}, \eqref{eq:Slater50}, and \eqref{eq:Slater51}, respectively.


\begin{theorem}[Initial values for the $\mathcal S_1$ basis]\label{thm:S1-initial}
The initial values for $B_{\mathcal S_1}(n,h)$ are
\begin{equation*}
    B_{\mathcal S_1}(0,h)=\begin{cases}
        1 &  \text{if $h=0$,}\\
        0 &  \text{otherwise.}
    \end{cases}
\end{equation*}
\begin{equation*}
    B_{\mathcal S_1}(1,h)=\begin{cases}
        q &  \text{if $h=1$,}\\
        0 &  \text{otherwise.}
            \end{cases}
\end{equation*}
\end{theorem}
\begin{theorem}[Basis recurrence for $\mathcal S_1$]\label{thm:S1-recurrence}
For any $n,h\geq1$, $B_{\mathcal S_1}(n,h)$ satisfies the following recurrence.
\begin{equation}\label{eq:S1Rec}
B_{\mathcal S_1}(n,h)=zq^{2h}B_{\mathcal S_1}(n-2,h-1)+azq^{2h}B_{\mathcal S_1}(n-2,h-2).
\end{equation}
\end{theorem}

\begin{proof}
Given a partition $\lambda$ in $\mathcal{B}_{\mathcal{S}_1}$ with length $n$ and largest part $h$, then the second largest part of $\lambda$ can be either $h$ or $\overline{h}$. In the first case, the third largest part of $\lambda$ will be $h-1$ and the weight of the first two parts is $zq^{2h}$, while in the second case, the third largest part of $\lambda$ is $h-2$, and the weight of the first two parts is $azq^{2h}$. So, by deleting the first two parts from $\lambda$, we have the desired recurrence.
\end{proof}


\begin{theorem}[Closed basis forms for $\mathcal S_1$]\label{thm:S1-closed}
The closed forms for $B_{\mathcal S_1}(n,h)$ are given by
\begin{equation}\label{eq:S1B1}
B_{\mathcal S_1}(2n,n+h)=a^hz^nq^{n(n+1)+h(h+1)}{n\brack h}_{q^2},
\end{equation}
and
\begin{equation}\label{eq:S1B2}
B_{\mathcal S_1}(2n+1,n+h+1)=a^hz^nq^{(n+1)^2+n+h(h+1)}{n\brack h}_{q^2}.
\end{equation}
\end{theorem}

\begin{proof}
We start with \eqref{eq:S1B1}, note that
\begin{align*}
B_{\mathcal S_1}(2n,n+h)=&a^hz^nq^{n(n+1)+h(h+1)}{n\brack h}_{q^2}\\
=&a^hz^nq^{n(n+1)+h(h+1)}\left(q^{2h}{n-1\brack h}_{q^2}+{n-1\brack h-1}_{q^2}\right)\\
=&a^hz^nq^{n(n+1)+h(h+1)+2h}{n-1\brack h}_{q^2}+a^hz^nq^{n(n+1)+h(h+1)}{n-1\brack h-1}_{q^2}\\
=&zq^{2n+2h}\cdot a^hz^{n-1}q^{n(n-1)+h(h+1)}{n-1\brack h}_{q^2}\\
&+azq^{2n+2h}\cdot a^{h-1}z^{n-1}q^{n(n-1)+h(h-1)}{n-1\brack h-1}_{q^2}\\
=&zq^{2n+2h}B_{\mathcal S_1}(2n-2,n+h-1)+zaq^{2n+2h}B_{\mathcal S_1}(2n-2,n+h-2),
\end{align*}
so \eqref{eq:S1B1} satisfies \eqref{eq:S1Rec}. Similarly, for \eqref{eq:S1B2},
\begin{align*}
B_{\mathcal S_1}(2n+1,n+h+1)=&a^hz^nq^{(n+1)^2+n+h(h+1)}{n\brack h}_{q^2}\\
=&a^hz^nq^{(n+1)^2+n+h(h+1)}\left(q^{2h}{n-1\brack h}_{q^2}+{n-1\brack h-1}_{q^2}\right)\\
=&a^hz^nq^{(n+1)^2+n+h(h+1)+2h}{n-1\brack h}_{q^2}+a^hz^nq^{(n+1)^2+n+h(h+1)}{n-1\brack h-1}_{q^2}\\
=&zq^{2n+2h+2}\cdot a^hz^{n-1}q^{n^2+n+h(h+1)-1}{n-1\brack h}_{q^2}\\
&+azq^{2n+2h+2}\cdot a^{h-1}z^{n-1}q^{n^2+n+h(h-1)-1}{n-1\brack h-1}_{q^2}\\
=&zq^{2n+2h+2}B_{\mathcal S_1}(2n-1,n+h)+zaq^{2n+2h+2}B_{\mathcal S_1}(2n-1,n+h-1).
\end{align*}
So, they both satisfy the desired recurrence. By checking the initial values, we finish the proof.
\end{proof}


\end{document}
